\chapter*{Abstract}
\addcontentsline{toc}{chapter}{Abstract}

Conventional data-driven approaches to transit planning frequently predict where
transit \emph{already exists} rather than where it is \emph{needed}, producing
circular models that reward the status quo. This thesis develops a demand-driven,
transferable framework that instead diagnoses \emph{unmet} transit demand and
generates candidate corridors to address it, applied to the Guadalajara
Metropolitan Area (\zmg; \num{1881} census enumeration areas, \ageb, across ten
municipalities).

The framework proceeds in eight linked stages: (i)~a Tier-1 transit-demand surface
estimated from census, economic, and street-network data via a doubly-constrained
gravity model; (ii)~calibration of the distance-decay parameter against the 2022
Origin--Destination survey; (iii)~a supply layer measuring \textsc{gtfs}-based
accessibility and an independent \emph{coverage-gap} index that becomes the
supervised target, removing the circularity of ``has-a-stop'' labels; (iv)~a
Node--Place--People prioritization map with an explicit equity term; (v)~a formal
multi-objective evaluation function; (vi)~a corridor generator based on frontier
anchors, a minimum-spanning-tree diameter trunk, and an anchor-directness
feasibility gate; (vii)~an audit of the existing network; and (viii)~a validation
suite combining hold-out backtests, benchmarking, and before/after coverage
metrics.

The diagnostic layer is the strong contribution: the coverage-gap index is
corroborated out-of-sample by the 2025 opening of \zmg Line~4, whose corridor shows
\num{1.6} times the metropolitan high-gap rate. The generative layer produces at
least one substantive, feasible, merit-passing corridor while honestly
characterizing its limitations. Finally, the entire framework is transferred
end-to-end to two additional metropolitan areas of contrasting size---Toluca and
Aguascalientes---demonstrating that it is genuinely portable rather than bespoke to
\zmg.

\vspace{1em}
\noindent\textbf{Keywords:} public transit; travel demand; coverage gap; Node--Place
model; multi-objective optimization; equity; \textsc{gtfs}; transferability;
Guadalajara.

\clearpage
